\section{Conclusion}\label{sec:conclusion}

We have assembled a fully reproducible pricing pipeline for a hybrid
real-estate investment structure inspired by the synthetic CDO mechanism
but applied to the rental cash flows of a Paris residential building.
Every parameter is calibrated on publicly available French data; every
numerical result reported here is traceable to the same Monte Carlo run
that wrote \texttt{artifacts/results.csv}.

The headline finding is that the structure is feasible: a senior tranche
attaching at 60\% prices close to par with a probability of negative
annualised return of \probNegSenior{} on the Paris-intermediate
calibration, while the equity tranche delivers a long-tailed payoff with
\probNegEquity{} probability of loss. The choice of credit model matters
quantitatively: switching from the Gaussian one-factor copula to a
Student-t copula with $\nu = 5$ drops the senior fair-value-to-par from
\seniorFairGauss{} to \seniorFairStudent{} and widens the fair coupon by
roughly twelve basis points, the empirical analogue of the
Donnelly-Embrechts critique. The two classical benchmarks --- single-
apartment ownership and the SAS pool --- carry distinct risk profiles
even when matched on expected return, so the tranching choice is not
cosmetic.

The pipeline is open source and available at
\texttt{tranche-pricing-paris}.
